\section{Method}

\subsection{Overview}

HGRN2-GatedDeltaNet replaces standard attention with a **two-branch linear 
recurrent module**. Both branches process the sequence sequentially with a 
fixed-size state, giving $\mathcal{O}(L)$ complexity. A learned gate fuses 
the two branches per token.

\textbf{Figure 4.1} shows the pipeline.

\begin{lstlisting}
    Input X $\in$ R^{B$$\times$L$$\times$d}
              |
       +------+------+
                      
    +------+    +----------+
    | HGRN |    | DeltaNet |
    | -2   |    |          |
    +--+---+    +-----+----+
       |              |
       |  $\gamma$ = $\sigma$(W_g X)|
       |     |        |
                       
       +-------------------+
       | Y = $\gamma$ Y$_1$ + (1-$\gamma$) Y$_2$|
       +---------+---------+
                 |
                  
              out_proj
                 |
                  
              Output Y
\end{lstlisting}

\subsection{HGRN-2 Branch}

The HGRN-2 branch (Qin et al., 2024) computes:

\[
u_t = \text{SiLU}(W_u x_t + b_u)
\tag{4.1}
\]

\[
\alpha_t = \sigma(W_\alpha x_t + b_\alpha)
\tag{4.2}
\]

\[
s_t = \alpha_t \odot s_{t-1} + (1 - \alpha_t) \odot u_t
\tag{4.3}
\]

The output is $y_t^\text{HGRN} = s_t$ (or optionally projected).

\textbf{Initialization.} $s_0 = 0$. Gates $W_\alpha$ are initialized so that 
$\alpha_t \approx 0.5$ initially.

\subsection{Gated DeltaNet Branch}

The Gated DeltaNet branch maintains a matrix state per head:

\[
S_t^h = \alpha_t S_{t-1}^h + \beta_t (v_t^h - S_{t-1}^h k_t^h) (k_t^h)^\top
\tag{4.4}
\]

where:

\begin{itemize}
\item $k_t^h = W_k x_t / \|W_k x_t\|$ (normalized key).
\item $v_t^h = W_v x_t$ (value).
\item $\alpha_t, \beta_t \in (0, 1)$ are scalar gates.

\end{itemize}
The output is:

\[
y_t^{h, \text{Delta}} = S_t^h q_t^h
\tag{4.5}
\]

where $q_t^h = W_q x_t$.

\textbf{Why the Delta rule?} The update $\beta_t (v_t - S_{t-1} k_t)$ writes only 
the \textbf{error} into memory, avoiding redundant writes.

\subsection{Fusion Gate}

The fusion gate is a linear projection:

\[
\gamma_t = \sigma(W_g x_t + b_g)
\tag{4.6}
\]

with $W_g \in \mathbb{R}^{d \times d}$ and $b_g \in \mathbb{R}^d$.

\textbf{Initialization.} We set $W_g = 0$ and $b_g = 0$, giving 
$\gamma_t = 0.5$ initially (equal blend). During training, the gate learns 
to prefer the branch that best fits each token.

\textbf{Why the gate matters.} At initialization, both branches are 
approximately equally weighted. If one branch's \texttt{copy\_weights} fails, we 
bias the gate toward the surviving branch (see  4.6).

\subsection{Combined Output}

The final output is:

\[
Y = \gamma \odot Y^\text{HGRN} + (1 - \gamma) \odot Y^\text{Delta}
\tag{4.7}
\]

followed by an output projection:

\[
\text{out} = W_o Y + b_o
\tag{4.8}
\]

\subsection{Failure Handling}

\subsubsection{Both Branches Fail}

If both \texttt{copy\_weights} calls fail, we \textbf{abort} with \texttt{RuntimeError}. This 
prevents injecting an untrained random module.

\subsubsection{One Branch Fails}

If only HGRN succeeds:
\begin{lstlisting}
self.gate.weight.zero_()
self.gate.bias.fill_(4.0)  # sigmoid(4) $\approx$ 0.98
\end{lstlisting}

If only DeltaNet succeeds:
\begin{lstlisting}
self.gate.weight.zero_()
self.gate.bias.fill_(-4.0)  # sigmoid(-4) $\approx$ 0.02
\end{lstlisting}

This ensures the gate leans toward the surviving branch.

\subsubsection{Output Projection Fails}

If \texttt{out\_proj} weights cannot be found, we \textbf{abort} with \texttt{RuntimeError}. 
This prevents injecting a random output projection.

\subsection{Algorithm}

\textbf{Algorithm 1:} HGRN2-GatedDeltaNet forward pass
\begin{lstlisting}
Input:  X $\in$ R^{B$$\times$L$$\times$d}, attention_mask
Output: Y $\in$ R^{B$$\times$L$$\times$d}

 1: Y_hgrn $\leftarrow$ HGRN2Branch(X)                   Sequential O(L)
 2: Y_delta $\leftarrow$ DeltaNetBranch(X)               Sequential O(L)
 3: $\gamma$ $\leftarrow$ $\sigma$(W_g X + b_g)                        (B, L, d)
 4: Y $\leftarrow$ $\gamma$   Y_hgrn + (1-$\gamma$)   Y_delta
 5: Y $\leftarrow$ out_proj(Y)
 6: return Y
\end{lstlisting}

\subsection{Complexity Analysis}

\begin{table}[H]
\centering
\small
\begin{tabular}{lll}
\toprule
\textbf{Component} & \textbf{Time} & \textbf{Memory} \\
\midrule
HGRN-2 & $\mathcal{O}(L d)$ & $\mathcal{O}(d)$ \\
DeltaNet & $\mathcal{O}(L H d^2)$ & $\mathcal{O}(H d^2)$ \\
Gate & $\mathcal{O}(L d^2)$ & $\mathcal{O}(d)$ \\
out\_proj & $\mathcal{O}(L d^2)$ & $\mathcal{O}(d^2)$ \\
\textbf{Total} & $\mathcal{O}(L H d^2)$ & $\mathcal{O}(H d^2)$ \\
\bottomrule
\end{tabular}
\end{table}

Compare to standard attention: $\mathcal{O}(L^2 H d)$. Our method is 
\textbf{linear} in $L$.

\subsection{Parameter Analysis}

\textbf{HGRN2-GatedDeltaNet does NOT reduce parameter count.}

\begin{table}[H]
\centering
\small
\begin{tabular}{lll}
\toprule
\textbf{Component} & \textbf{Baseline} & \textbf{Ours} \\
\midrule
Q projection & $d^2$ & $d^2$ (frozen) \\
K projection & $d^2$ & $d^2$ (frozen) \\
V projection & $d^2$ & $d^2$ (frozen) \\
O projection & $d^2$ & $d^2$ (frozen) \\
HGRN-2 params & 0 & $\sim 4d^2$ \\
DeltaNet params & 0 & $\sim 4d^2$ \\
Gate & 0 & $d^2$ \\
LoRA & 0 & $8 r_L d$ \\
\bottomrule
\end{tabular}
\end{table}

\textbf{Total increase:} ~9$d^2$ per layer.

For $d = 768$:

\begin{table}[H]
\centering
\small
\begin{tabular}{ll}
\toprule
\textbf{Config} & \textbf{Params per layer} \\
\midrule
Baseline & 2.36M \\
HGRN2-GatedDeltaNet & ~7.7M \\
\textbf{Increase} & \textbf{+226\%} \\
\bottomrule
\end{tabular}
\end{table}

This is a significant parameter increase. We report it honestly.

\subsection{Training}

We train only:
\begin{itemize}
\item The gate $W_g, b_g$.
\item LoRA adapters on Q, K, V, O, and the branch-specific projections.

\end{itemize}
All base weights from the pretrained attention and the two branches are 
frozen.

\textbf{Trainable params:} for $d = 768$, $r_L = 16$:
\begin{itemize}
\item LoRA: $8 \cdot 16 \cdot 768 = 98,304$
\item Gate: $768^2 + 768 = 590,592$
\item \textbf{Total:} ~689K per layer

\end{itemize}
For 12 layers: ~8.3M trainable.
\begin{lstlisting}

---
\end{lstlisting}